% Consolidated textbook-ch12 from 1699b8435d575b3f9a89b4f32b3cd9072025e92c; source fragments preserved.
% 2026-11-24 八批严格复核第4批：旧片段哈希对应历史版本，当前含二校修订。

% === BEGIN textbook-ch12-choices.tex blob a6bd54e7f4dc7548980b3c2746fc64040272963c ===
% 扫描《大学物理简明教程》第5版第12章量子物理基础
% 选择12.1.1—12.1.14，原书/PDF316—317页。原创推导，原题题意简述。
\subsection{选择题：量子辐射、光子、氢谱与量子态}

\sourceproblem{12.1.1}{316}{316}
\problemtext 绝对黑体应具备怎样的吸收/反射性质？
\workings
绝对黑体按定义吸收入射的\emph{全部波长}电磁辐射，
无反射、无透射，故题中“不能反射任何光线”正确，选\textbf{D}。
\crosscheck “黑体”不等于在任意温度下不辐射可见光；
热黑体仍会按温度发射热辐射。

\sourceproblem{12.1.2}{316}{316}
\problemtext 钠的光电红限波长$540$ nm，某单色光使最大光电子动能为$1.2$ eV；
求入射光的波长。
\workings
逸出功$W=hc/(540\,\mathrm{nm})$，
爱因斯坦方程$hc/\lambda=W+1.2$ eV，
采用$hc\simeq1240\,\mathrm{eV\,nm}$，
\[
\boxed{\lambda=\frac{1240}{1240/540+1.2}\,\mathrm{nm}
\simeq354.6\,\mathrm{nm}}.
\]
选\textbf{D（355 nm）}。
\crosscheck 入射光要比红限更短波，结果满足$\lambda<540$ nm。

\sourceproblem{12.1.3}{316}{316}
\problemtext 将地球与金星均视为理想黑体，温度分别$293$ K和$773$ K，
且两者尺寸近似相同，比较单位表面面积的辐射出射度。
\workings
斯特藩--玻尔兹曼律$M=\sigma T^4$，故
\[
\boxed{\frac{M_{\rm Venus}}{M_{\rm Earth}}
=\left(\frac{773}{293}\right)^4\simeq48.44}.
\]
选\textbf{D（48.4）}。
\crosscheck 两者\emph{单位面积}比值与球体半径无关；对总辐射功率则还需面积。

\sourceproblem{12.1.4}{316}{316}
\problemtext 同种金属的红限频率$\nu_0$，以频率$\nu_1,\nu_2$照射，
第二次遏止电压是第一次的两倍；求$\nu_2$与$\nu_1,\nu_0$的关系。
\workings
$eU_a=h(\nu-\nu_0)$，令$U_{a,2}=2U_{a,1}$，
得$\nu_2-\nu_0=2(\nu_1-\nu_0)$，故
\[
\boxed{\nu_2=2\nu_1-\nu_0}.
\]
选\textbf{C}。
\crosscheck 这不是频率直接加倍，因逸出功引入频率阈值。

\sourceproblem{12.1.5}{316}{316}
\problemtext 康普顿散射中散射后波长为入射波长的1.2倍；
求散射光子能量与反冲电子动能之比。
\workings
光子能量反比波长，$E'_{\gamma}=E_\gamma/1.2=5E_\gamma/6$；
电子原先静止时反冲动能$K=E_\gamma-E'_\gamma=E_\gamma/6$。
\[
\boxed{E'_\gamma/K=5}.
\]
选\textbf{D}。
\crosscheck 忽略外界束缚能；能量守恒适用于自由电子的康普顿散射。

\sourceproblem{12.1.6}{316}{316}
\problemtext 氢原子巴耳末线系中波长最长与次长的谱线分别记$\lambda_1,\lambda_2$，
求两者比值。
\workings
最长线$n=3\to2$，次长线$n=4\to2$，
由$1/\lambda=R_H(1/2^2-1/n^2)$，
\[
\boxed{\frac{\lambda_1}{\lambda_2}
=\frac{1/4-1/16}{1/4-1/9}=\frac{27}{20}}.
\]
选\textbf{C}。
\crosscheck 最大波长应对应能级差最小的巴耳末跃迁。

\sourceproblem{12.1.7}{317}{317}
\problemtext 静止氢原子由$n=3$直接跃迁到$n=1$，
原子质量$1.67\times10^{-27}$kg，估算反冲速度。
\workings
光子能量$\Delta E=13.6(1-1/9)=12.09$ eV，
动量$p_\gamma\simeq\Delta E/c$；
原子反冲动量大小相等，
\[
\boxed{v_r\simeq\frac{12.09(1.602\times10^{-19})}
{(1.67\times10^{-27})(3.00\times10^8)}
\simeq3.87\,\mathrm{m/s}}.
\]
选\textbf{A（4 m/s）}。
\crosscheck 原子反冲能远小于跃迁光子能量，忽略它只产生极小修正。

\sourceproblem{12.1.8}{317}{317}
\problemtext 氢原子由能级$E_3$直接到$E_1$发出$\lambda_3$，
分两步$E_3\to E_2\to E_1$发出$\lambda_2,\lambda_1$；
求波长关系。
\workings
在忽略原子反冲、外场及精细结构并比较同一组三个能级时，由能级差相加
$hc/\lambda_3=hc/\lambda_2+hc/\lambda_1$，
所以
\[
\boxed{\lambda_3^{-1}=\lambda_1^{-1}+\lambda_2^{-1}}.
\]
选\textbf{C}。
\crosscheck 光子\emph{频率/波数}可相加，波长不可直接相加。

\sourceproblem{12.1.9}{317}{317}
\problemtext 在德布罗意与戴维孙--革末实验相关的四个论断中，
判断哪些真实反映微观粒子的波粒二象性。
\workings
光子具有波粒二象性，实物粒子也具有波粒二象性；
戴维孙--革末电子衍射实验验证了德布罗意物质波。
“电子只有粒子性”错误。故正确为命题1、2、3，选\textbf{B}。
\crosscheck 不可用经典轨道的粒子图像排除电子的干涉与衍射。

\sourceproblem{12.1.10}{317}{317}
\problemtext 电子由静止通过电势差$U$加速后德布罗意波长为$0.04$ nm，
求加速电压。
\workings
非相对论$p=\sqrt{2m_e eU}$，
$\lambda=h/p$；用常数式
$\lambda(\mathrm{nm})\simeq1.226/\sqrt{U(\mathrm V)}$，
\[
\boxed{U\simeq(1.226/0.04)^2\,\mathrm V
\simeq940\,\mathrm V}.
\]
选\textbf{D}。
\crosscheck 940 eV对511 keV的电子静能较小，非相对论近似合理。

\sourceproblem{12.1.11}{317}{317}
\problemtext 两种不同静质量的粒子具有相同德布罗意波长，
比较它们的动量、速度、动能与动能比。
\workings
由$\lambda=h/p$，两种粒子的\emph{动量大小}严格相同：
\[
\boxed{p_1=p_2=h/\lambda}.
\]
选\textbf{A}。
\crosscheck 同一动量下非相对论动能$p^2/(2m)$一般不同，
速度也一般不同。

\sourceproblem{12.1.12}{317}{317}
\problemtext 势阱$-a\le x\le a$内给归一化波函数
$\psi(x)=a^{-1/2}\cos(3\pi x/2a)$，求$x=5a/6$处概率密度。
\workings
$|\psi(x)|^2=a^{-1}\cos^2(3\pi x/2a)$，
代入$x=5a/6$，角度为$5\pi/4$，故
\[
\boxed{\rho(5a/6)=\frac1a\cos^2\frac{5\pi}{4}=\frac1{2a}}.
\]
选\textbf{A}。
\crosscheck 概率密度量纲是$\mathrm{length}^{-1}$，不是无量纲概率。

\sourceproblem{12.1.13}{317}{317}
\problemtext 判断不确定性原理是否禁止单独确定坐标、单独确定动量，
以及它是否适用于所有微观粒子。
\workings
不确定性原理限制的是同一态中位置与动量分量的不确定度\emph{同时}减小，
$\Delta x\,\Delta p_x\ge\hbar/2$；
单独某一量可以任意提高测量或制备精度；理想位置、动量本征态是广义本征态，不等同于可归一化且两者方差都有限的物理态。
因此正确的是命题(c)、(d)，选\textbf{B}。
\crosscheck 原题纸面把$\hbar$写成$h$样式；
采用标准量子力学常量$\hbar=h/(2\pi)$阐释，不将不同常量混用。

\sourceproblem{12.1.14}{317}{317}
\problemtext 电子处在原子K壳层（$n=1$），四种$(n,l,m_l,m_s)$组合中何者允许？
\workings
主量子数$n=1$时只能$l=0$，相应$m_l=0$，
自旋投影$m_s=\pm1/2$；
因此
\[
\boxed{(n,l,m_l,m_s)=(1,0,0,-1/2)}
\]
是选项\textbf{A}所列的允许态。
\crosscheck 源题选项B写$l=0,m_l=-1$，不满足$|m_l|\le l$；
C、D不属于$n=1$壳层。
% === END textbook-ch12-choices.tex ===

% === BEGIN textbook-ch12-blanks.tex blob 5a66b4775f00d2fd14f36fbf725985cf48c0b157 ===
% 量子物理基础填空题12.2.1—12.2.13；用户原书/PDF第317—318页。
\subsection{填空题：光电截止、氢谱与测不准关系}

\sourceproblem{12.2.1}{317}{317}
\problemtext 波长300 nm照射某金属，光电子动能由0到$4.0\times10^{-19}$ J，
求遏止电压及红限频率。
\workings
由$eU_a=K_{\max}$，有
\[
\boxed{U_a=\frac{4.0\times10^{-19}}{1.602176634\times10^{-19}}
\approx2.50\,\mathrm V}.
\]
光子能$hc/\lambda\approx6.62\times10^{-19}$ J，
逸出功$W=hc/\lambda-K_{\max}$，
\[
\boxed{\nu_0=W/h
=c/\lambda-K_{\max}/h\approx3.96\times10^{14}\,\mathrm{Hz}}.
\]
\crosscheck 红限频率必须为正且低于入射光频率$10^{15}$ Hz。

\sourceproblem{12.2.2}{317}{317}
\problemtext X射线分别在$45^\circ$及$60^\circ$发生康普顿散射，
比较散射波长增量$\Delta\lambda_1:\Delta\lambda_2$。
\workings
由$\Delta\lambda=\lambda_C(1-\cos\theta)$，
\[
\boxed{\frac{\Delta\lambda_1}{\Delta\lambda_2}
=\frac{1-\cos45^\circ}{1-\cos60^\circ}
=2-\sqrt2\approx0.5858}.
\]
\crosscheck 60度散射增量更大，所得比值小于1。

\sourceproblem{12.2.3}{317--318}{317--318}
\problemtext 氢原子电子由能量$-0.85$ eV跃迁到$-3.4$ eV；
求释放光子能量及两状态主量子数。
\workings
$E_n=-13.6/n^2$ eV，$-0.85$ eV对应$n=4$，
$-3.4$ eV对应$n=2$，
\[
\boxed{E_\gamma=3.4-0.85=2.55\,\mathrm{eV}},
\qquad \boxed{n=4\to2}.
\]
\crosscheck 波长约$1240/2.55=486$ nm，是巴耳末可见谱线。

\sourceproblem{12.2.4}{318}{318}
\problemtext 依据原图$n=1,2,3,4$的能级及箭头，
选出所发光子波长最短、频率最低的跃迁。
\workings
波长最短对应能级差最大，故
$\boxed{n=4\to n=1}$；
频率最低对应所画向下跃迁中能级差最小，故
$\boxed{n=4\to n=3}$。
\crosscheck 这里比较的是图中\emph{发射箭头}，并非任意能级差或吸收跃迁。

\sourceproblem{12.2.5}{318}{318}
\problemtext 氢原子基态能量$E_1=-13.6$ eV，求$n=4$能量。
\workings
\[
\boxed{E_4=-\frac{13.6}{4^2}=-0.85\,\mathrm{eV}}.
\]
\crosscheck 相对电离极限零能，束缚态能量均为负。

\sourceproblem{12.2.6}{318}{318}
\problemtext 某电子在原子激发态中平均停留$10^{-8}$s；
估算能级能量最小不确定度。
\workings
若采用教材的能时数量级估计$\Delta E\,\Delta t\gtrsim\hbar/2$并把寿命作为特征时间尺度，
\[
\boxed{\Delta E\gtrsim\frac{\hbar}{2(10^{-8})}
=5.27\times10^{-27}\,\mathrm J
\approx3.29\times10^{-8}\,\mathrm{eV}}.
\]
\crosscheck 能时不确定度不同于位置--动量算符的严格同类形式，
自然寿命线宽（稳定末态、指数衰减、无其他展宽）的能量全高半宽为$\Gamma=\hbar/\tau\approx6.58\times10^{-8}$ eV。洛伦兹谱的方差并非有限，故不能把该全宽与严格标准差下界混同，也不能仅凭寿命宣称一个普适的最小能量标准差。

\sourceproblem{12.2.7}{318}{318}
\problemtext 用波长$\lambda$表达光子的能量、动量大小与相对论质量等效值。
\workings
\[
\boxed{\varepsilon=\frac{hc}{\lambda}},\qquad
\boxed{p=\frac h\lambda},\qquad
\boxed{m_{\rm eq}=\frac{\varepsilon}{c^2}=\frac h{\lambda c}}.
\]
\crosscheck 光子的\emph{静质量为零}；$m_{\rm eq}$仅为原书旧术语下的能量等效质量。

\sourceproblem{12.2.8}{318}{318}
\problemtext 氢原子电子主量子数$n=3$，求允许轨道角动量大小最大值。
\workings
$l=0,1,2$，最大$l=n-1=2$，
\[
\boxed{L_{\max}=\sqrt{l(l+1)}\hbar=\sqrt6\,\hbar}.
\]
\crosscheck 不用玻尔旧量子论的$L=n\hbar$取代现代量子角动量。

\sourceproblem{12.2.9}{318}{318}
\problemtext 电子限于宽度$\Delta x=5$ nm的区域，
按\emph{原题指定的估算关系}$\Delta x\Delta p_x\ge h$，
求动量分量的最小估计不确定度。
\workings
依题面特定估算式
\[
\boxed{\Delta p_x\gtrsim\frac h{\Delta x}
=\frac{6.62607015\times10^{-34}}{5\times10^{-9}}
\approx1.33\times10^{-25}\,\mathrm{kg\,m/s}}.
\]
\crosscheck \textbf{源式与标准严格界限不同：}标准海森堡不等式为
$\Delta x\Delta p_x\ge\hbar/2=h/(4\pi)$。
若$5$ nm明确为位置标准差，标准不等式给$\Delta p_x\ge1.06\times10^{-26}$ kg\,m/s；若仅表示支持区间全宽$L=5$ nm，则$\sigma_x\le L/2$，可推出至少$\sigma_p\ge\hbar/L=2.11\times10^{-26}$ kg\,m/s，不能将区间宽度自动当作标准差。
本条为忠于\emph{原题指定公式}而计算，不能混称两式等价。

\sourceproblem{12.2.10}{318}{318}
\problemtext 用波函数$\Psi(\bm r,t)$描述微观粒子，
写出$\Psi\Psi^*$的概率解释、波函数要求及归一化条件。
\workings
\[
\boxed{\rho(\bm r,t)=\Psi\Psi^*=|\Psi|^2},\qquad
\boxed{\int_{\mathbb R^3}|\Psi(\bm r,t)|^2\,d^3r=1}.
\]
对于通常单粒子束缚态，波函数需平方可积、满足哈密顿量定义域及相应边界条件；常规有限势下常取连续可微代表，奇异势或无限势阱边界则按具体匹配条件处理。平面波等散射广义本征态采用不同的归一化约定。
\crosscheck 若只描述一维粒子，则归一化体积分应改为$\int|\psi(x)|^2dx$。

\sourceproblem{12.2.11}{318}{318}
\problemtext 电子自旋磁量子数$m_s$可能的两个值。
\workings
\[
\boxed{m_s=+\frac12}\qquad\boxed{m_s=-\frac12}.
\]
\crosscheck 电子自旋量子数为$s=1/2$，而$m_s$为其投影量子数。

\sourceproblem{12.2.12}{318}{318}
\problemtext 主量子数$n=2$的一个电子壳层最多有多少种单电子量子态？
\workings
计入自旋的态数为$2n^2$，
\[
\boxed{N=2(2)^2=8}.
\]
\crosscheck $n=2$由$2s$（2个态）和$2p$（6个态）组成。

\sourceproblem{12.2.13}{318}{318}
\problemtext 多电子原子中电子排布遵从哪两项基本原则？
\workings
依原书填\textbf{泡利不相容原理}与\textbf{能量最小原理}。
\crosscheck 实际多电子组态还需结合洪特规则等判断简并亚层占据方式；
本题只要求两项教材概括原则。
% === END textbook-ch12-blanks.tex ===

% === BEGIN textbook-ch12-solutions-001-009.tex blob 2862cda691cc616f72bc73ead5cd70b3ad8958ea ===
% 《大学物理简明教程》第5版量子物理基础：解答12.3.1—12.3.9。
% 原书/PDF第318—319页；题意独立摘要，带核验的定量推导。
\subsection{解答题：黑体、光电效应与康普顿散射}

\sourceproblem{12.3.1}{318}{318}
\problemtext 把三颗恒星当理想黑体，辐射谱的峰值波长分别为
$0.55,0.35,0.29\,\mu$m，求对应表面温度。
\workings
对按单位波长定义的黑体谱，维恩位移定律$\lambda_mT=b$，
$b=2.89777\times10^{-3}\,\mathrm{m\,K}$，得
\[
\boxed{\begin{aligned}
T_1&=b/(0.55\,\mu\mathrm m)\simeq5.27\times10^3\,\mathrm K,\\
T_2&=b/(0.35\,\mu\mathrm m)\simeq8.28\times10^3\,\mathrm K,\\
T_3&=b/(0.29\,\mu\mathrm m)\simeq9.99\times10^3\,\mathrm K.
\end{aligned}}
\]
\crosscheck 峰值波长越短，黑体温度越高；此处是黑体\emph{色温}近似而非真实恒星大气模型。

\sourceproblem{12.3.2}{318}{318}
\problemtext 炉壁小孔的总辐射出射度为$22.8\,\mathrm{W/cm^2}$，
视作黑体辐射，求炉内温度。
\workings
$M=22.8\times10^4\,\mathrm{W/m^2}$。
由$M=\sigma T^4$，
$\sigma=5.670374419\times10^{-8}\,\mathrm{W/(m^2K^4)}$，
\[
\boxed{T=\left(\frac{22.8\times10^4}{\sigma}\right)^{1/4}
\simeq1.416\times10^3\,\mathrm K}.
\]
\crosscheck 注意$1\,\mathrm{W/cm^2}=10^4\,\mathrm{W/m^2}$。

\sourceproblem{12.3.3}{318}{318}
\problemtext 铝金属逸出功$4.2$ eV，以波长$200$ nm光照射，
求最大动能、遏止电压和红限波长。
\workings
光子能量$E=hc/\lambda\simeq1240/200=6.20$ eV，
光电效应给
\[
\boxed{K_{\max}=E-W=2.00\,\mathrm{eV}},\qquad
\boxed{U_a=K_{\max}/e=2.00\,\mathrm V}.
\]
红限
\[
\boxed{\lambda_0=hc/W\simeq(1240/4.2)\,\mathrm{nm}
\simeq295\,\mathrm{nm}}.
\]
\crosscheck 入射光波长200nm小于红限295nm，能发生外光电效应。

\sourceproblem{12.3.4}{318}{318}
\problemtext 波长$300$ nm紫外光照射某金属，
最大光电子速度$5.0\times10^5$ m/s，求金属红限波长。
\workings
非相对论电子动能
$K_{\max}=\tfrac12m_ev^2\simeq1.14\times10^{-19}$ J
$\simeq0.711$ eV。
入射光子能量约$1240/300=4.133$ eV，
则逸出功$W\simeq3.422$ eV，故
\[
\boxed{\lambda_0\simeq\frac{1240}{3.422}\,\mathrm{nm}
\simeq362\,\mathrm{nm}}.
\]
\crosscheck $\lambda_0>300$ nm是光电逸出发生的必要条件。

\sourceproblem{12.3.5}{318}{318}
\problemtext 原图12.3.5给钠金属遏止电压随入射频率的拟合直线：
$\nu_0\simeq4.39\times10^{14}$ Hz，且$(\nu,U_a)$约有
$(5.5\times10^{14}\,\mathrm{Hz},0.47\,\mathrm V)$及
$(9.6\times10^{14}\,\mathrm{Hz},2.13\,\mathrm V)$。
由图估计截止频率、普朗克常量及逸出功。
\workings
截止频率是图的横轴截距，
\[
\boxed{\nu_0\simeq4.39\times10^{14}\,\mathrm{Hz}}.
\]
从光电方程$eU_a=h\nu-W$知斜率为$h/e$。
按图上两个数值点估算
\[
\boxed{h_{\rm graph}
=e\frac{2.13-0.47}{(9.6-5.5)\times10^{14}}
\simeq6.49\times10^{-34}\,\mathrm{J\,s}}.
\]
逸出功约为
\[
\boxed{W\simeq h_{\rm graph}\nu_0
\simeq2.85\times10^{-19}\,\mathrm J
\simeq1.78\,\mathrm{eV}}.
\]
\crosscheck 这些是\emph{图读估算值}，
不能把由粗线和圆点读出的$h$当成精确定义的$6.62607015\times10^{-34}$；
不同两点选取可能略有偏差。特别地，给出的两个近似点所确定直线的横轴截距为$4.339\times10^{14}$ Hz，不恰等于单独图读的$4.39\times10^{14}$ Hz。若严格使用两点直线，应统一采用其截距，得到$W\approx2.81\times10^{-19}$ J（约1.76 eV）；待核原图并统一拟合。

\sourceproblem{12.3.6}{319}{319}
\problemtext 人眼在理想暗适应条件下每秒吸收5个波长500nm的蓝绿色光子
刚能形成可见刺激，求每秒光子携带的能量及对应入眼功率。
\workings
单光子能量$\varepsilon=hc/\lambda$，
每秒5光子所携能量
\[
\boxed{E_{\rm 1s}=5\frac{hc}{500\,\mathrm{nm}}
\simeq1.99\times10^{-18}\,\mathrm J}.
\]
相应被吸收的光功率（入眼功率还取决于吸收效率）
\[
\boxed{P=E_{\rm 1s}/(1\,\mathrm s)
\simeq1.99\times10^{-18}\,\mathrm W}.
\]
\crosscheck 辐射实际感知阈值依观察时间、视网膜吸收效率和刺激部位而变；
本题只是在理想吸收的光子计数模型下换算能量；若吸收效率为$\eta$，入射功率为上述吸收功率除以$\eta$。

\sourceproblem{12.3.7}{319}{319}
\problemtext 设一光子的能量等于电子静止能量，
求光子的频率、波长与动量。
\workings
由$\varepsilon=m_ec^2$，
\[
\boxed{\nu=\frac{m_ec^2}{h}
\simeq1.236\times10^{20}\,\mathrm{Hz}},
\]
\[
\boxed{\lambda=\frac{h}{m_ec}
\simeq2.426\times10^{-12}\,\mathrm m},\qquad
\boxed{p=\frac{\varepsilon}{c}=m_ec
\simeq2.731\times10^{-22}\,\mathrm{kg\,m/s}}.
\]
\crosscheck 此波长等于\emph{电子康普顿波长}（非约化），不是电子本身的静止德布罗意波长。

\sourceproblem{12.3.8}{319}{319}
\problemtext 波长$0.0708$ nm的X射线在石蜡中发生康普顿散射，
求散射角$90^\circ$和$180^\circ$时的散射波长。
\workings
$\lambda'=\lambda_0+\lambda_C(1-\cos\theta)$，
电子康普顿波长$\lambda_C=h/(m_ec)=0.00242631$ nm。
故
\[
\boxed{\lambda'_{90^\circ}\simeq0.073226\,\mathrm{nm}},
\qquad
\boxed{\lambda'_{180^\circ}\simeq0.075653\,\mathrm{nm}}.
\]
\crosscheck 后向散射产生最大波长增加$2\lambda_C$。

\sourceproblem{12.3.9}{319}{319}
\problemtext 康普顿实验中波长$0.015$ nm的光子被原先静止的自由电子散射，
在$120^\circ$方向求散射波长相对增量。
\workings
\[
\Delta\lambda=\lambda_C(1-\cos120^\circ)
=\frac32\lambda_C\simeq0.0036395\,\mathrm{nm}.
\]
故
\[
\boxed{\frac{\Delta\lambda}{\lambda_0}
=\frac{0.0036395}{0.015}\simeq0.2426},
\qquad
\boxed{\lambda'\simeq0.01864\,\mathrm{nm}}.
\]
\crosscheck 题问\emph{波长变化量与原波长之比}，
不是$\lambda'/\lambda_0\simeq1.2426$，也不是光子能量比。
% === END textbook-ch12-solutions-001-009.tex ===

% === BEGIN textbook-ch12-solutions-010-018.tex blob df3e555b87ecf93d7f1342b1246a068fd15a0229 ===
% 《大学物理简明教程》第5版量子物理基础计算12.3.10—12.3.18
% 原书/PDF第319页，独立推导；自然线宽与光子位置问题说明模型定义。
\subsection{解答题：氢谱、物质波、线宽和量子态计数}

\sourceproblem{12.3.10}{319}{319}
\problemtext 基态氢原子吸收能量为$12.75$ eV的光子。
求激发后的能级及返回低能级过程中可能发出的谱线。
\workings
氢原子$E_n=-13.6/n^2$ eV，吸收后
$E_f=-13.6+12.75=-0.85$ eV，对应
\[
\boxed{n=4}.
\]
不考虑额外选择定则限制、只按能级能差枚举的可能发射跃迁有
\[
\boxed{4\to3,\quad4\to2,\quad4\to1,\quad
3\to2,\quad3\to1,\quad2\to1}.
\]
共有$4(4-1)/2=6$种不同能量差。
\crosscheck \emph{严格原子跃迁}还要求电偶极选择定则（如$\Delta l=\pm1$），
源题未提供初态亚能级，故上述是仅依据主量子数能级的\emph{可选谱线集合}，
不是断言同一个原子一次退激发同时发出六个光子。

\sourceproblem{12.3.11}{319}{319}
\problemtext 基态氢原子受单色光激发后，
其辐射谱中仅出现两条不同的巴耳末谱线。
求这两条谱线波长及激发光频率。
\workings
巴耳末系终态为$n=2$。
只有$n=3\to2$和$n=4\to2$两条，
说明激发的最高主量子数按题目能级模型取$n_{\max}=4$。
\[
\boxed{\lambda_{3\to2}
=\frac1{R_H(1/2^2-1/3^2)}
\simeq656.5\,\mathrm{nm}},
\]
\[
\boxed{\lambda_{4\to2}
=\frac1{R_H(1/2^2-1/4^2)}
\simeq486.3\,\mathrm{nm}}.
\]
由$1\to4$的激发能$\Delta E=12.75$ eV，
\[
\boxed{\nu_{\rm exc}=\frac{12.75\,e}{h}
\simeq3.08\times10^{15}\,\mathrm{Hz}}.
\]
\crosscheck 两条谱线的数目依赖激发后的级联通道可能遍历；
采用玻尔/黎德伯主能级理想模型，忽略细结构、选择定则和实际跃迁概率；以上用氢原子黎德伯常量$R_H\approx1.09678\times10^7\,\mathrm{m^{-1}}$计算真空波长。常用空气谱线值约656.3和486.1 nm，不能混称为同一介质波长。若不假定所有主能级通道可见，仅凭观察到两条谱线不能严格唯一反推出最高激发能级。

\sourceproblem{12.3.12}{319}{319}
\problemtext 基态氢原子吸收$12.09$ eV光子后，
电子玻尔轨道半径增加了多少倍？
\workings
$13.6(1-1/n^2)\simeq12.09$ eV，解得$n=3$。
玻尔半径$r_n=n^2a_0$，
\[
\boxed{r_{\rm after}=9r_{\rm before}}.
\]
若“增加多少倍”指增量与原值之比，
\[
\boxed{\frac{\Delta r}{r_1}=9-1=8}.
\]
\crosscheck 原半径变为\emph{9倍}、增加了原半径的\emph{8倍}，
两种中文表述不能混用。

\sourceproblem{12.3.13}{319}{319}
\problemtext 一光子和一个电子的波长都为$0.20$ nm，
求它们各自的动量大小与总能量。
\workings
两者按德布罗意关系动量大小相同：
\[
\boxed{p_\gamma=p_e=\frac{h}{\lambda}
=\frac{6.62607015\times10^{-34}}{0.20\times10^{-9}}
\simeq3.313\times10^{-24}\,\mathrm{kg\,m/s}}.
\]
光子的能量
\[
\boxed{E_\gamma=pc=\frac{hc}{\lambda}
\simeq9.93\times10^{-16}\,\mathrm J
\simeq6.20\,\mathrm{keV}}.
\]
电子\emph{总能量}必须包含静能：
\[
\boxed{E_{e,\rm total}=\sqrt{p^2c^2+m_e^2c^4}
\simeq8.19\times10^{-14}\,\mathrm J
\simeq511.0\,\mathrm{keV}}.
\]
\crosscheck 与电子总能量不同，其动能仅约
$p^2/(2m_e)\simeq37.6$ eV。源题特别问\emph{总能量}，
不能把37.6 eV误当总能量。

\sourceproblem{12.3.14}{319}{319}
\problemtext 自由电子动能$200$ eV，求德布罗意波长；
再判断它与直径为1nm的孔或障碍物相遇时波性是否值得考虑。
\workings
因为$K\ll m_ec^2=511$ keV，
\[
\boxed{\lambda=\frac{h}{\sqrt{2m_eK}}
\simeq0.0867\,\mathrm{nm}}.
\]
与尺寸$d=1.0$ nm的孔或障碍物相比，
$\lambda/d\simeq0.087$，
衍射的典型特征角数量级
$\theta\sim\lambda/d\simeq0.087$ rad。
所以该电子\textbf{既具有粒子性也具有波动性}，
若观测角分辨足够精细，波动衍射效应不能先验忽略。
\crosscheck 不能只因$\lambda<d$就断定“没有波动性”；
实际可否测出衍射需考虑相干性、几何与探测分辨率。

\sourceproblem{12.3.15}{319}{319}
\problemtext 某激发态回到基态辐射谱线中心波长约400nm，
观察到线宽$\Delta\lambda=10^{-5}$nm，
估算原子激发态的平均寿命。
\workings
窄谱线近似下
\[
\Delta\nu\simeq\frac{c\,\Delta\lambda}{\lambda^2}
=\frac{(3.00\times10^8)(10^{-14})}{(4.00\times10^{-7})^2}
\simeq1.875\times10^7\,\mathrm{Hz}.
\]
若线宽指\emph{自然洛伦兹线型的频率全高半宽}，且末态稳定、无多普勒展宽等额外展宽，
量子态概率寿命$\tau$与线宽满足
$\Delta\nu_{\rm FWHM}=1/(2\pi\tau)$，因此
\[
\boxed{\tau\simeq\frac1{2\pi\Delta\nu}
\simeq8.49\times10^{-9}\,\mathrm s}.
\]
\crosscheck \textbf{宽度定义影响常数系数}：
若教材采用粗略能时估计$\Delta E\,\Delta t\sim h$
则得到约$5.33\times10^{-8}$s；
两者的差异来自$2\pi$以及谱宽/寿命的定义，
不可称为同一个严格数值。

\sourceproblem{12.3.16}{319}{319}
\problemtext 一光子的波长为300nm，波长测量相对精度为百万分之一，
估算其纵向位置的不确定尺度。
\workings
波长误差
$\Delta\lambda/\lambda=10^{-6}$，故
$\Delta p\simeq(h/\lambda)(\Delta\lambda/\lambda)$。
只有额外假定给出的相对精度表征波包本身的波长标准差而非仪器测量误差，才可将下式作为标准位置--动量不确定性下界的数量级估计，
\[
\Delta x\gtrsim\frac{\hbar}{2\Delta p}
=\frac{\lambda}{4\pi(\Delta\lambda/\lambda)}
=\boxed{2.39\times10^{-2}\,\mathrm m}.
\]
\crosscheck \textbf{光子位置算符有理论细节}，
本式是沿传播方向的单光子波包长度的启发性估计，
不是任何光子位置测量的普适严格算符下界；单独的仪器测量精度不能决定光子波包的内禀位置分布。
按其他教材宽度约定（如$\Delta x\sim\lambda/\delta$）
结果可在数倍或$4\pi$因子上不同。

\sourceproblem{12.3.17}{319}{319}
\problemtext 一维无限深势阱$0\le x\le a$内粒子波函数为
$\psi(x)=\sqrt{2/a}\sin(3\pi x/a)$；
求概率密度达到最大的位置。
\workings
\[
\rho(x)=|\psi(x)|^2=\frac2a\sin^2\frac{3\pi x}{a}.
\]
最大值$2/a$要求
$3\pi x/a=(k+\tfrac12)\pi$，
即$x=(2k+1)a/6$。
在$[0,a]$内取$k=0,1,2$，
\[
\boxed{x=a/6,\quad a/2,\quad 5a/6}.
\]
\crosscheck 端点$x=0,a$概率密度为零，正好满足无限深势阱边界条件。

\sourceproblem{12.3.18}{319}{319}
\problemtext 原子电子以$n,l,m_l,m_s$四量子数表征，
分别在(1)固定$n,l,m_l$，
(2)固定$n,l$，(3)仅固定$n$时，求可有的不同量子态数。
\workings
固定$n,l,m_l$时电子自旋投影$\pm1/2$，
\[
\boxed{N_1=2}.
\]
固定$n,l$时$m_l=-l,\ldots,l$共$2l+1$种，
各有两种自旋，故
\[
\boxed{N_2=2(2l+1)}.
\]
固定$n$时对$l=0,\ldots,n-1$求和，
\[
\boxed{N_3=2\sum_{l=0}^{n-1}(2l+1)=2n^2}.
\]
\crosscheck 这些是\emph{不同单电子态}的最大计数；
若是\emph{同一原子实际电子占据数}还须遵守泡利不相容原理。
% === END textbook-ch12-solutions-010-018.tex ===
